\documentclass{ltjsbook}
\usepackage{docmute} 
\usepackage{../../myStyle} 

\begin{document}
\chapter{社会心理学のコア}
\label{sp01_introduction}

\section{ご挨拶}

それでは時間になりましたので，授業の方を始めさせていただきたいと思います。私はこの大学に来るのは初めてで，不慣れなものですから，システムとかよくわかってないところもあります。授業の進め方として皆さんから何かできればリアルタイムに反応が欲しいなと思うんですけど，そういう良いツールがとりあえず思いつかなかったので，こういうのを用意してあります。slido.comというところがあるんですが，ここにアクセスしていただきますと，Q\&Aみたいなのが書き込めるようになってますので，何かコメントなどがあったら，無記名で書き込めますので，なんでも反応をお願いします。もちろん記名で質問とか，「今のところもう少し詳しく」みたいなのがあったら書いてください。対応します。本当はこちらのLMSを使った方がいいのかなと思ってるんですが，よくわかりません。とりあえず授業の終わりに先ほどアンケート機能を使って「授業のコメントください」というのを書いてございますので，それにもご協力いただければなと思います。

さて，今日は1回目の授業，初回の授業ということですので，社会心理学の特殊講義と，社会心理学とは何かというあたりをお話しさせていただくんですが，内容についてというか，僕は何を喋ろうとしているかというような話をしたいと思うんです。で，言葉が難しいんですが，ちょっと私，今回この授業はですね，古い大学の授業スタイルでやらせてもらおうかなと思っています。スライドとか使わずに，濤濤と喋る感じにしようかなと思っているんです。よっぽどのことがあれば，投影資料とか動画とか持ってくるかなと思うんですけど。

ちょっと本筋からそれるんですけど，統計の授業をやっている時は教えなきゃいけない内容というのが割とカチッと決まっているので，僕はテキストをちゃんと作るんです。1年間だとか半期間の間で「ここからここまでのことを教えないといけないので，この順番で教えよう」と設計するんですけれども，この授業に関しては，もちろん一応大まかなアウトラインが決めてあるんですが，事前にその喋る内容を全部ガチガチに決めてしまわないでおこうという方針でいくつもりです。むしろ皆さんの方で，僕が喋ることをうまくまとめていただいて，皆さんなりに何というのかな，「小杉が言いたかったこと」，あるいは「社会心理学についての皆さんなりの答え」が，脳となり頭の中なりでできていただくということを目標にしようかなと思っているので，事前にこのキーワードとこれととかっていうのをガチガチに固めるつもりなく登壇しております。

この授業でスライドを作らないのもそれが理由で。スライドを作ってしまうと，資料として提示する点では分かりやすいんですけど，そういう授業のスライドっていうのは，得てして喋る側の人間の備忘録でしかないんですね。「次，これ喋らないと」「次，これ喋らなきゃ」みたいなことを書いてあるだけになったりするので，そういう，何というのかな，過剰書き的なとか，情報としての授業みたいなのは別の機会に譲っておいて，僕としては，なんかもうちょっと「古いスタイルの大学っぽい」みたいな言い方しましたけど，にゅるっとですね，全体像が出てくるような授業にしたいなと思っていますので，最初からテキストあるいは配布資料を作るというつもりはございません。

その代わりといってはなんですが，今，ちょっと片方だけイヤホンを入れてあるんですけど，ボイスメモを録ってありまして。あとでこれから文字起こしをします。多少整頓して，授業録みたいな形で資料として公開できればなというふうに思っています。その時に，ちゃんと引用文献リストを付けるからね。

古いやり方のこの授業のデメリットは，引用文献リストとか出典みたいなのが，記憶だとか，そういうことで喋ってしまうと間違っている可能性があるということです。書いた文字にするとちゃんとその辺が是正されるという良さがあるんですが，長短所それぞれ授業のやり方，スタイルがありますので，その勢いに任せて喋って，それで引用とかが喋ってた時に間違ってたようなことがあったら，後で補正するという形でいこうと思っていますので，ご承知の上，聞いてください。


さて，と。何を言うべきだったかな，最初に。

\section{自己紹介をさせてください}

まず，ちょっと自己紹介がてら話の中に入っていきたいと思うんですが，私，社会心理学の授業をするのがすごい久しぶりなんです。私，1976年の生まれで，大学に入ったのが94年ですから，もう20年ぐらい前の話なんです。「お前の年齢なんか興味ない」と思うかもしれませんけど，ちょっと本編とも絡んでくるので，ちょっと私の来歴の話をさせてくださいね。

94年に大学に入って，大学院まで6年，O.D.1年あたり，7年ぐらい大学にいたんですが，入ったのは関西学院大学っていう関西の私大です。で，社会心理学の教室があって，大学院があってというところに進んだわけですが，社会学部の中にある心理学コースでした。正確に言うと，社会学部の中に理論専攻というのがあって，その中に数理社会学と社会心理学，何か知らんけど心理学がそこに入ってるというところでして，そこで学位を取りましたんで，私の学位は社会学なんです。

私の学位は社会学なんですけど，やってることは指導教員が社会心理学でしたから，バリバリの社会心理学なんです。やってたことはそうなんですけど，その後，仕事を得まして山口大学という，山口県にある国立大学に行きました。行ったところは教育学部だったんですね。教育学部の中に心理学専攻がございまして，そこで授業を担当することになりました。担当する授業が統計と社会心理学の2つです。社会心理学とグループダイナミックス及び心理統計を教えること，というふうになってました。

今でもそうですが，そういうところは結構あるんじゃないかなと思います。何かというと，心理学のコースに統計の専門家がいることってなかなかなんですね。社会心理学の先生が統計を教えてるとか，実験の先生が統計を教えてるということがあって，僕もそのタイプで，社会心理学と統計を教えていました。社会心理学をやってたんですけど，今言ったように院生時代は隣に数理社会学がいて，副査に数理社会学の先生がいましたし，数学的な勉強仲間もいたというようなこともあって，あと僕も数学的なことがちょっと好きだったんで，統計を教えるのは得意な方でした。そのおかげで，統計の授業を担当するようになったんですね。それが山口大学時代で，主にこの2種類の授業をやってたわけですけれども，その後こっち来ました。7～8年目になりますから，2017年から18年頃に専修大学の方に異動してきました\footnote{正確には2018年4月から専修大学でした。}。

この専修大学のポストというのが心理統計だったんです。ベイズ統計とか教える心理統計枠がありまして，心理統計専門の枠に入りました。専修大学には社会心理学の先生が別にいらっしゃいまして，今でも忘れませんが，入社の面接の時にですね，「うちには社会心理学の先生がいるから社会心理学は教えなくていいよ」と言われたんですよ。「そうですか」ということで，それで僕は今はだからずっと心理統計ばかりを本職として教えているんです。

今回，こちらからお声がけをいただきまして，社会心理学を喋ってと言われたので，個人的には身に余る光栄というか，嬉しくて仕方がないんですけれども，さて，皆さんの方が不安になっているんじゃないですかね。「本当にお前が社会心理学を教えられるのか？」と。僕も，「それでいいんだろうか」とちょっと思ったりするんですけれども。
ということで，皆さんに考えていただきたいんですけど，私は社会心理学者を名乗ってよいでしょうか。

言い換えると，「社会心理学者とは，何をしていれば社会心理学者と言えるのか？」というようなことを考えてみたい，あるいは考えてもらいたいです。一応私は社会心理学会には入っております。社会心理学系の学会といえば，社会心理学会とグループダイナミックス学会というのがあります。歴史的にはグループダイナミックス学会の方が古いですね。僕，グループダイナミックス学会にも入ってたんです。ですが，専修に移るときに，いろんなところに顔を出しすぎているなというのもあったので，グループダイナミックス学会は辞めさせてもらいました。今は社会心理学会だけに所属しています。あ，その他にもちろん，日本心理学会とか行動計量学会とかには入っているんですが。

この「社会心理学会に入ってたら社会心理学者なの？」というところなんですけど，どうでしょうかね。なんでこんな話をするかというと，何が社会心理学で何をやってたら社会心理学者なのか，というところから，この授業は考えていきたいからです。

社会心理学会が50周年だったときにシンポジウムがありまして，早稲田の竹村先生という方が，社会心理学ってどんな学問かというのを50年間振り返るというシンポジウムをされました\footnote{2009年10月10-12日，大阪大学で開催された日本社会心理学会第50回大会・日本グループダイナミックス学会第56回大会合同大会における，記念シンポジウム「新たな社会心理学の展開と現状からの脱却」での話。}。その時に言ってはったのが，「社会心理学者というのは，社会心理学研究に論文を載せているから，社会心理学の研究をしている人だ」という話です。社会心理学会が出している『社会心理学研究』というところに，社会心理学に関する研究が載っているんですけど，そこで引用している論文というのが社会心理学系の文献を当然引用しているわけですね。これは当たり前の話なんですが，ただ，ちょっと見ていただいたら後で皆さんに調べてもらう時間を作ろうかなと思うんですけど，社会心理学系の論文というのは，結構自分の中の分野の引用が多いんです。例えば医学系の分野を引用するとか，工学の話を引用するとかっていうのはそんなに多くないんです。社会心理学という分野では，国内外問わず，社会心理学の論文を引用し合うことが多い。社会心理学会に入って，社会心理学の論文を引用して，社会心理学の論文を書いたら，社会心理学者なんじゃないか，という話になってしまうわけです。ちょっと極端な例ではありますが。それって本当に社会心理学者としての定義になるんでしょうか？という疑問を持ってしまいます。

何をもって社会心理学と言えるのか，その定義が非常に曖昧だという点で，これは社会心理学の特異なところでもあるし，学問としての難しさでもあります。どんな人でも，そこに論文を投稿して，それが受理されれば社会心理学者だと言われるのであれば，ある意味オープンな学問であるとも言えるかもしれません。しかし，私がこの授業で皆さんにお話ししたいのは，そういう「論文を書いて投稿すれば社会心理学者だ」ということではなく，社会心理学とは何か，社会心理学を学ぶということがどういうことなのか，というより本質的な部分です。

\section{社会心理学のゴールはどこにあるのか}

さて，社会心理学というのはどういう学問なのか？何を研究していたら，それは社会心理学と呼ばれるのか？社会心理学のゴールはどこにあるのか？そういうことを考えていくことが，この授業のテーマのひとつです。最初に言っておきますが，僕がこの授業で引用する論文やテキストは，世界最先端のホットトピックスというわけではありません。むしろ少し古い論文や，古典的な理論や本の話をすることが多いかもしれません。ですので，もし皆さんが今の最先端の社会心理学の研究を期待しているのであれば，もしかするとあまり役に立たないかもしれません。

でも，もし皆さんが長く社会心理学に関わることがあれば，あるいは心理学全般に関わることがあれば，社会心理学という学問体系がどういう流れで成り立ち，今後どうなっていくべきかという，メタ的な視点で捉えることが重要だと思います。僕はそのメタ的な視点を皆さんに提供することで，社会心理学という学問を少しでも理解する一助になれればと思っています。

では，全体の流れに戻りますが，まず社会心理学とは何か？というのがこの授業の最終的なテーマです。皆さん，もし「社会心理学って何？」と聞かれたらどう答えるでしょうか？多分，入門書を読むのが最初の一歩でしょう。幸い，社会心理学の教科書や入門書はたくさんありますし，書店でも「社会心理学」というタイトルのテキストが多く見つかります。それらを見ていただければ，社会心理学がどういう研究をしているのかは，大まかに分かると思います。

私が今日持ってきたのは，確かに少し古いのですが，2016年に出版された『社会心理学概論』という本です\parencite{SPgairon2016}。ナカニシヤ出版から出た本で，これもいわゆる概論です。この本は久しぶりに分厚くて細かい字がびっしり詰まった骨太なテキストでしたので，久しぶりにこういうテキストが出たなと思って買ったのを覚えています。内容的には，社会心理学の古典的な研究から始まり，さまざまなトピックが書かれています。

具体的には，第1章が社会心理学の古典的研究というタイトルで，これまでどんな研究がなされてきたのかをまとめています。その後も，2章が社会的認知，3章ステレオタイプ，4章社会的推論，5章自己過程と自己制御，6章感情と道徳，7章態度と説得，8章対人行動，9章人間関係，10章コミュニケーション，11章組織と集団家庭，12章集団間関係，13章インターネット，14章が文化，15章進化的アプローチ，16章アクションリサーチ，17章が健康，18章が環境，19章が規範と法，そして第20章が社会神経科学，となっています。全部で20章からなるこの本の内容を見ていると，社会心理学が取り扱うトピックが非常に幅広いことがわかるでしょう。

別に全部を覚える必要はありませんが，この章立てを見ただけでも，社会心理学がどれだけ多岐にわたる分野を含んでいるかが感じられると思います。
では，ちょっと質問ですが，今の章タイトルを聞いて，社会心理学がどういう研究をしているかというイメージが湧きますか？ どうですか？他の学問分野と比べて，少し作りが違うと思いませんか？

社会心理学の教科書というのは，だいたいこんな感じなんです。他にも，僕が授業をするにあたっていくつか昔のテキストを見ましたが，どれも同じように，いくつかのトピックスについてツラツラと書かれています。例えば，僕のかつての師匠が書いたテキストなんかは，「チャートでわかる社会心理学」というもので\parencite{SPchart1994}，図式化されて人間の行動のパターンが描かれています。「こういう状況ではこういう行動を取る」といった，過剰書きのように書かれている箇所があるんです。

しかし，今述べたように多くのジャンルが網羅的に書かれているからといって，それがそのまま社会心理学者の定義になるわけではありません。それぞれのトピックを知っているからといって，社会心理学者だと言えるかは別問題です。では，一体何を持って社会心理学者だと言えるのか？ここで少し皆さんに実感してもらいたいと思います。

ちょっとした試みですが，皆さん，パソコンやスマホをお持ちなら「社会心理学研究」という論文をJ-STAGEで検索してみてください。J-STAGEというのは，日本語論文の集積サイトでして，そこに社会心理学会が発行している論文集『社会心理学研究』があります。そこに様々な研究が掲載されていますので，興味がある論文のタイトルをいくつか見つけて，アブストラクト（要約）を読んでみてください。過去数年分遡って，ザラッと見ていただければいいかなと思います。そして，もし何か興味を惹かれる論文があったら，そのタイトルをコピーして，こちらにシェアしてもらえますか？

5分から10分くらい時間をとりますので，皆さん，J-STAGEで「社会心理学研究」を探してみてください。「巻号一覧」っていうのを見ると，40巻ありますから，たくさんいろんな研究がズラッとでてきますね。興味があるタイトルやアブストラクトを読んでみて，面白そうなものがあればシェアしていただければと思います。では，少しお待ちしますね。

……（少し待つ時間）……


では，ちょっと話をこちらに戻らせていただきます。はい，いろいろ挙げていただきました。ここで，皆さんに考えてもらいたいのです。これ，「社会心理学研究」っていう雑誌の中から皆さん選んでもらいましたね。これが社会心理学の研究として，社会心理学会が出している社会心理学者たちに読まれて，共有されている情報なわけです。これをザッと見たときに，もし社会心理学研究から取ってきましたというようなヒントがなければ，これ，本当に「社会心理学の研究だ」というふうに皆さん，受け取ってもらえるでしょうか？

これ，例えば子供のなんとかっていう話がありましたね。子供の虚言行動になんとかっていうのは\footnote{\textcite{KJ00003725146}と\textcite{Asano20232022-006}を混ぜてはなしちゃってます。ごめんなさい。}，人格心理学っぽくないですか？ 子供用セルフコントロール尺度の作成\parencite{Asano20232022-006}とか，子供を相手にして尺度を使って何かを測りましょうっていうことですね。これらから社会心理学って，連想できますか？観光写真調査法による観光地の魅力評価\parencite{Hayashi20191819}。心理学かどうかもちょっとタイトルだけ見た感じだと，よくわからないですよね。僕，この研究をやってる人は友達だから，あんまり無碍に，悪し様に言いたいわけではないですが，これが社会心理学の王道かって言われると，土真ん中かって言われると，なんか違うような気もしますね。

分配の正義とリスク化の意思決定\parencite{Kotani20211914}。政治学みたいな話なような気もしますし，あるいは実験経済学みたいな話のような気もする。文化差ですとか，大学生における，とかっていうと，一部の現代的な何かを世相を反映したものかな，というような気がしますが，この「高レベル放射性廃棄物地層処分施設」の社会的受容とくると\parencite{Shibata20242214}，これは社会学，あるいは社会運動の話かなのような内容かと思いますよね。「愛着不安」とかというのは\parencite{Kanemasa20171618}，臨床心理学的なテーマなような気もします。「反移民的態度」とかっていうのも\parencite{Kashihara20222014}，政治的な話のような気もしますが，まあまあまあ，ことほどさように，非常に多岐にわたっているのが，今の社会心理学なわけです。

先ほどこの「社会心理学概論」ですが，というようなテキストもご案内して，そこの章立ての話をしましたけれども，それもこれも似たような感じで，小さなトピックスがいっぱい網羅的にあるような世界になっているって思いませんか。

で，例えば社会心理学者として登壇して，授業をしますよ，皆さん授業してください，とかって言われたら，皆さんが社会心理学を研究してたら，何をやりましょうか。
一つ一つのこの話のジャンルって，もちろん同じ雑誌に載ってるんですから，同じカテゴリーなような気もしますけど，それにしても，例えばこの「異性愛者のジェンダー自尊心と同性の同性愛者に対する態度」っていう研究\parencite{Suzuki2015842}をずっとやっておられた研究者がですよ，社会心理学の授業全15回をやりましょうって言ったときに，15回も全部これの話しますかね。きっとしないでしょうね。

社会心理学っていうのは全体はこれぐらいでねっていうような話をされると思うんです。もちろんご自身の専門に関するところを深く掘り下げる。それの最新トピックスを紹介するっていうのが大学の授業としていい授業だと思うんですが，僕がそれこそ山口大学に赴任して，さあ，社会心理学の授業をしなさいと言われたときに困ったのは，今言ったようなことです。社会心理学の授業として何を教えたらいいんだろうかです。

こういった類の教科書があって，各ジャンル，社会的認知，個人内過程，対人的な対人行動，小集団の話，みたいなトピックスはもちろんいっぱいあるんです。一応このジャンルに，この学会に身を置いて，人の話を聞いてきましたから，私も一応それぞれこういう研究をしてるのだろうなというような当たりはつきます。でも，こういった各論がいっぱいあるんですけど，社会心理学とはこういった学問であるという総論といいますか，社会心理学の原論みたいなのが掴めないんですよね。

それが，この社会心理学という領域の難しいところで，さて，社会心理学を教えましょうねとなったときに，何を教えたら社会心理学を教えたことになるのか。あるいは皆さんが学んだときには社会心理学ってこうだったんだねって残るものは何か，ということを考えてもらいたいんです。

私の経験談ですが，昔，高校生の時の日本史の先生が割と研究者肌の先生で，しかもそれも弥生時代ぐらいのところを大学で専門にやっておられたらしい人がいました。日本史の授業ですけど，夏休みが来てもまだ平安ぐらいまでしか行ってない。1年間で絶対日本の歴史を終わらせないといけないんじゃないの，というようなことがありました。でも歴史をやってるなという，コアのところは掴めたように思うんですね。

そんな感じで，例えば，こういったどれかのトピックをぐーっと深くやったとして，社会心理学とはこういうものなのだな，という社会心理学のコアに当たる部分が見つかるでしょうか。見つからないのではないかというのが，私の印象です。皆さんも同じように感じてくれてるかな，と思うのですが。それぞれの研究者がそれぞれの研究テーマでもって，内輪で集まってやってるっていうのはいいんですが，さて，社会心理学の本質はどこか，というようなことが見えないんじゃないかなと思うわけです。

できればどうにも私はこの授業の中で，僕の考える社会心理学っていうのはここがコアなんじゃないですか，というような話をしたいと思っています。僕が中途半端で，社会心理学に専念してないというか，自分が社会心理学者なのかどうなのか，自分の定義はよくわからないんですが，まあでも，社会心理学を中心ではないにしろずっと考えてきた人間が，何について考えたらいいのか，ということについてお話しできればなあと思っています。　

この授業は，おそらく皆さんが普段聞いたことのないような話をすることになるかなと思います。方法論的にマニアックな話をしようと思っていますし，今言った学問全体の考え方みたいなところを中心にお話ししようと思っていますので，繰り返しになりますが，ホットトピックスみたいなのはパラパラとは出てきませんけれども，そういったようなことが皆さんと一緒に考えていけたらいいかなというふうに思っています。

\section{一般的な導入}

さてここで，社会心理学の概略，あるいは全体像に関わるような一般的な導入をしておきたいと思います。

皆さんどこかで聞いたことがあるような話かもしれませんが，そもそも社会心理学における話のややこしい問題のひとつに，「社会」というのは何か，という話があります。Social Psychologyですから，「社会」というのは「Social」，Societyの訳だというのはおわかりいただけるかと思いますけれども，そこから入ってみたいと思います。

これは私がかつて，社会学の授業を聞いてたときに，なるほどうまいこと言うなと思ってしみじみ感じ入ったところなんですが。「社会」というのは何かという問いについて，社会学は「社会」というものがあるよ，という前提から話が始まります。ここで「社会」というのは何か，というのは中心的な問題なんですけど，「社会とは何か」というようなことを言ったときに，その社会学のテキスト\parencite{SPkiso1990}では，まず``Society''という言葉がいつ「社会」と翻訳されるようになったか，というところから話を始めたんですね。

そこで引用していたのが，岩波文庫の『翻訳語成立事情』\parencite{SPhonyaku}という岩波新書ですね。これ，面白いので皆さんも是非読んでみてください。要は，明治になって文明開化したときに外国からの言葉が入ってくるわけですね。それを日本語に翻訳しないといけないということになったときに，みんな困りましたよという話なんです。この本には，いろいろ翻訳に困ったキーワードが10個挙げられてるんですが，そのうちの一番が「社会」です。

「社会」というのは外国語で``Society''で，それを翻訳するときに福沢諭吉さんとかが，これどう訳したらいいのか，と困ったという話が書いてあります。福沢諭吉は，``Society''を日本語に翻訳するときには，「交際」，人付き合いですね。「交際」とか「交わり」，つまり「人間交際」，「国」，「世人」という訳を当てていたそうです。

その他にも，いろんな訳語があって，中村正直という人は``Society''を訳して「政府」とか「仲間連中」とか，「人民の会社」といったりしてます。「会社」という言葉がわりと新しかったみたいですね。会社というのは，人がたくさん集まってチームを組んでいるみたいな意味ですが，「仲間」とか「人民の会社」，みたいな訳を``Society''に当てていた。政府という意味を当てていたこともあるそうです。

そういうところから，その後，「社会」という新しい言葉ができまして，そのおかげで，Society＝社会と訳すようになった。人がたくさんいて，みんなであれこれやり取りしていて，それでかつプライベートとはちょっと違う，何かパブリックな意味があるんだろうね，というような，人と交わるときのパブリックみたいな言葉の意味としてふわっとした，それまでにはなかった「社会」という語を当てたわけです。我々はもう訳すものだという癖がついていますけれども，明治のその頃になるまで，日本には「社会」という概念がそもそもなかったわけですね。

なので，おそらく社会心理学もその流れで来ているでしょうから，「社会心理学」と言われると，人間が関係していて，しかもそれがインタラクションしていて，ちょっとですね，パブリックな意味のある何か，というようなところが研究の対象になっていたんだろう，というふうに考えることができるかなと思います。

ちなみに，この「社会心理学」というのがスタートした頃は，社会学とかともあんまり大きな違いはなかったわけですし，社会心理学の周辺領域といえば，社会学だとか，人文学だとか，人類学だとか，民族学だとか，そういうふうな学問分野みたいなのがわりと周りに多かったわけです。今や，社会心理学はインターネットとか，あるいは人工知能みたいなのもそうかもしれませんが，情報系の科学なんかの深い関わりがある領域になっておりますけれども，昔はもっともっと，なんかこうベタベタの文系的な感じのものだったわけです。

ちなみにこの『訳語成立事情』の第2章が「個人」です。``Individual''。この``Individual''っていうのも，文明開化するまでは「個人」という概念が日本にはなかったんじゃないかという話ですね。他にも「近代」とか「美」とか「恋愛」とかもそうです。要は``love''みたいなのがどう訳したらいいかわからなかったんですよね。「お慕い申しております」みたいな感覚であって，「恋愛」というのが出てきたのが明治以降だということです。ほかにも「存在」とか「自然」とか「権利」「自由」「彼」「彼女」みたいなのも最初どう訳すべきか悩んだ，という話が書いてありますので，ちょっと興味があったら読んでみてください。

さて「社会」というのがこういう意味なんだよ，ということがあるんですけど，これを使っての社会心理学です。心理学という以上は心の要素が入ってくるわけですが，社会心理学の定義とは何かといいますと，皆さん，今僕が社会心理学の総体，コアみたいなのが無いですって僕が言いましたよね。そうすると「いやいや待てよ，学問として成立してるんだから，学会もあって心理学教室もあるんだから，学問の定義があるやろう」と思われるかもしれません。

一応，社会心理学業界で使われている社会心理学の有名な定義を読み上げておきます。Allport，G.W.という人がした定義\parencite{Allport1954}です。その人の定義というのがいわゆる古典的な社会心理学の定義としてよく使われています。読み上げますので，ちょっと聞いてください。社会心理学は「個人の思考，感情，行動が，どのようにして他の人間の現実の存在，あるいは想像上，ないし暗黙の存在によって影響を受けるのかを理解し，説明することを目指すものである」\footnote{原文も載せておきます。``The scientific investigation of how the thoughts, feelings and behaviours of individuals are influenced by the actual, imagined, or implied presence of others.''\parencite{Allport1954}}。
もう一回言います。「個人の思考，感情，行動が，どのようにして他の人間の現実の存在，あるいは想像上，ないし暗黙の存在によって影響を受けるのかを理解し，説明することを目指すものである」。

なるほど，と思われたでしょうか？僕は，実はこれを何回読んでも，何が言いたいのか，いまいちピンときません。まず，「個人の思考，感情，行動を理解し，説明することです」というのは，一応納得するとしても，具体的に言うと，思考だとか感情だとか行動というのは一応3つとも厳密には違うものですから，それを理解するためのものと言われても，一概に言えないという気はします。

そして，それがどうやって説明されるかというと，ここがややこしいポイントです。どうにかして「他の人間の現実の存在，あるいは想像上，ないし暗黙の存在によって影響を受けるのか」を理解するそうですが，その影響のソースが「他の人間の現実の存在，あるいは想像上，ないし暗黙の存在」と言っています。人の感情や行動が他の人の存在によって，僕がどう影響を受けるかを理解する，というのはわかりますが，「他の人がいるだけではなく，それが現実の存在，あるいは想像上の存在でも良い」というわけです。

つまり，いなくていいんですよ。極端な話，この教室に誰もいなかったとしても，僕が誰かがいるかもなっていうのを想像しながら動いていたら，それは社会心理学の研究テーマになるというわけです。想像上の存在でもいいし，暗黙の存在，つまり「いているような雰囲気を醸し出している」ぐらいでもいい。

それって確かに，僕の行動が誰もいない時と，誰かいるかもなーと薄々気づいている時と，誰かいるに違いないと想像しながらやっている時とで全部違うかもしれませんけど。でも，それって，そこに一人しかいないわけでしょ。それ，社会心理学？社会なの？人がいてますか？交わってますか？パブリックな要素ありますか？というようなことが気になりますよねえ。

多分オールポートは，非常に広い範囲のいくつかの研究を最大公約数的に説明できる何かを考えたんだとは思いますが，ちょっとですね，広げすぎなんじゃないか，と。特にこの「暗黙の存在」とか「想像上の存在」と言い出すところが面白い。

僕，昔，社会心理学の研究の被験者になったときに，「恋人との行動を研究したいので，恋人と普段どういうふうにやり取りしているかを答えてください。もし恋人があなたにいらっしゃらないのでしたら，いたと想像して答えてください」という調査を受けたことがあります。想像上の恋人と，実際の恋人ではだいぶ違うぞと個人的には思うんですけど，皆さんいかがですかね。それらを同じようなものとして扱って良いものか。少なくとも，社会心理学が扱おうとしているのは，物理的な他人だけではないということでしょうかね。みんなとか，全体とか，あるいは「他の人もいているかもな」と一般化された他者を想像した上での，そういうのも社会心理学の対象というのか，ということです。

今，Sli.doの方に面白いコメントが来まして，「幽霊を怖がれば社会心理学になるのか」，ということですが，この定義からいくと，社会心理学になってしまうっぽいです。それで良いのか？そこも含めて研究しないといけないのだ，というのだったらそれで良いのですが，少なくとも物理的な存在，あるいは客観的な行動，Behavioralな側面だけではアプローチしきれないものを社会心理学は扱おうとしている。それを扱うのこそ社会心理学なんだ，となるわけですね。これでいいと個人的には思ってますけど。

ただ少なくとも，それをやるんだ，という理屈が成立しないといけない。つまり社会心理学が考えているのは，物理的な存在だけではなく，実在的な存在を考えないといけないですね。共通認識としてみんながそこに何かがあると認めたのであれば，それはあるものとして話を進めましょう，ということを言っているからです。そうすると，サイエンスの基盤としていわゆる自然科学みたいなのを範にした科学の成り立ちとは，ちょっと違う仕組みで設計しないと研究っていうのができないんじゃないでしょうか。

\section{社会心理学の分類}

さて，ちょっと話が横にそれますが，公認心理師というのがございます。心理学界が作った資格でして，特に臨床系の資格として作られたものです。私の大学，専修大学も公認心理師という資格に対応した授業をやっているんです。この公認心理師という資格を出すためには，公認心理師資格課程に認められた授業を展開している必要があります。その授業の中では，社会心理学に関しては「社会・集団・家族心理学」といったものが当てられています。

この公認心理師科目というのは，なかなか手強い話なんです。科目名が完全に一致していないとダメだと言われるんです。科目名が一致していたら中身は別にどうでも，というわけでは無いと思いますが，少なくとも公認心理師に対応している授業科目名は「社会・集団・家族心理学」と指定されています。この3つがセットなのは間違いないです。社会心理学会は一応抵抗したんですよ。「社会心理学と集団心理学と家族心理学は全部別物です。社会心理学界には，あまり家族研究をやっている人がいないので，これで一括りにされても困ります」という意見を出したはずなんですけど，いろんな事情でこれがワンセットになりました。

これがワンセットになっているということは，公認心理師の人たちが受ける社会心理学とは，こういう社会心理学だということです。こういう分け方もあるんだ，というのはヒントとして知っておいてもらえてもいいかなと思います。確かに社会的な要素はあるんですけど，こうやってまとめられてしまうところにも，押し切られてしまったところにも社会心理学の限界があるんじゃないかなと思います。これこそが社会心理学のコアで，これを含んでないからこれは違うものです，と言えなかったんですね。

これも確かに社会の中身だよね。何か分からないから入れておこう，まぁまぁいいじゃない，みたいな感じでまとめられてしまっている，というのが現状ではないかなと思います。これも，各論がいっぱいあるという性質のせいですよね。

あ，そうそう，集団心理学というのも本当は間違いなんです。社会心理学の中でよく言われるいろんな分類の仕方があるんですが，一つは，社会心理学の分類の仕方として，サイズによる分類というのがあります。

社会心理学を大きく分けると，ミクロとメゾとマクロの3段階に分かれるといわれます。小・中・大なんですが，大体ミクロレベルの社会心理学というと，1人から3人までが対象ですね。まず1人の時点で社会心理学なのか，というのは気になりますが。1人ですから社会なのかと言うんですが，オールポートの定義的には，1人ぼっちでいたとしても，幽霊の存在とか誰かのことを思っていたら成立する，というのでいいのでしょう。

さっき紹介したテキスト\parencite{SPgairon2016,}の中にも「個人内過程」というジャンルがあったと思います。ことほど左様に，一人でも社会心理学というジャンルがあるんですが，1人から3人ぐらいのレベルの社会心理学のことをミクロ社会心理学としておきましょう。

次に，マクロというのはミクロの反対ですから，たくさんですね。これは流行現象ですとか，政治運動ですとか，そういった1人1人のレベルではなく，全体として出てくる現象についての社会心理学なんだ，と言われたりします。

注目はメゾ，真ん中らあたりというやつです。このメゾの心理学というのが，グループダイナミックスの領域になります。ダイナミックスという言い方をしたのは，なぜかというと，集団心理学ではなく，このジャンルは正確には「集団力学」だからです。

グループダイナミックスを作ったK.Lewinという人は，「ダイナミックス」のところにポイントがあったのです。これを分解して「グループサイコロジー」と言ってしまうのは，本来的な目的とは違う。なぜこの真ん中のところだけ取り出されているか，ということを強調しておきたいですね。

メゾのグループダイナミックスというのは，大体範囲としては3人，4人ぐらいから小中学校の1学級ぐらいのサイズだと思ってください。この教室にいらっしゃる皆さんぐらいのサイズです。つまり，顔と名前が完全に一致する人ばかりではないが，全く知らない中でもない，ぐらいの範囲であれば，グループダイナミックスが発動する世界になるかなと思います。僕も思い出してみると，小学校の時は3，40人学級でしたけど，全員の名前までは覚えていなかったのですが，この子は同じクラスの子，同じ2組の子だ，違うクラスの子だ，というぐらいは分かりましたからね。違うクラスに間違って入ってしまったら，違うクラス来たってのが分かるぐらいには全体感があったかなと思うので，それぐらいのレベルの話です。

ここが特にグループダイナミックスと言われて出てきているのはなぜかというと，しかもこれが社会心理学の一部なんですけど，なぜここだけ別の「学」がついてるかというと，おそらくこのミクロ的な，個人の心的なプロセスと，現象レベルで出てくる全体性みたいなものの混じり合うところが，一番大きな研究テーマになるからなんですね。社会と個人が混ざり合うところが，社会心理学の中心的な問題であるべきなんじゃないかと僕なんか思ったりするわけです。だから特にそこの部分だけ，ミクロとマクロの混じり合うダイナミックスというところを取り上げて，一つの学問ジャンルにした，というのが実質のところではないかなと思います。

社会心理学の研究をやっています，というと，大体どのぐらいのサイズの研究をやっていますか，ということで分かれたりはしますが，ミクロとマクロのダイナミックスというところをしっかり考えているかどうか，というのは大きな違いかなと思います。社会心理学者というのは，他の学問ジャンルに流れていくこともありまして，統計の方に流れた僕が言うのもおこがましいかもしれませんが，ミクロ社会心理学は，社会心理学は研究方法として実験とか調査とかしたりしますからね，そのこともあって最近では実験経済学とかとも仲良くなったりしています。

あるいは，実験経済学者がやっていることは，ほとんど社会心理学者が今までやってきたようなことじゃないか，とも言われています。TverskyとKahnemanがプロスペクト理論\footnote{プロスペクト理論（prospect theory）は，不確実性下における意思決定モデルの一つ。選択の結果得られる利益もしくは被る損害および，それら確率が既知の状況下において，人がどのような選択をするか記述するモデルをさす。(Wikipediaより)}でノーベル賞を取っていますけれども，実験刺激みたいなのを出した時に，人がどう反応するか，刺激の出し方によって変わります，という話は社会心理学では昔から言われていたことでもあるわけです。

そういう実験的なアプローチを社会心理学はやりますよ，と言った時に，実験心理学と社会心理学はどう違うんですか，認知心理学とどう違うんですか，ということは説明できないといけないですね。社会心理学も「社会的認知」という言葉があったりしますからね。あれがまた話がややこしいわけです。

最近は実験哲学というのも出てきましたね。哲学というジャンルは，実験というのだからほど遠いと思っていたんですけれども，実験哲学というのは，トロッコ問題みたいなやつを調査することによって，「みんなはどっち側にトロッコを進めさせた方が正しいと思っているか」を調査するってな研究の仕方をするらしいです。僕は実験哲学の本は1冊しか読んでませんから\footnote{ここで言及しているのは，\textcite{SPjikkentetugaku}です。}，それが全てじゃないよと言われたらごめんなさいと思うんですけど，少なくとも調査研究をすることによって，哲学的な答えの一つにしましょう，というのは，社会心理学とはどのように違うのか考えものですね。

これも社会心理学というのは何か，というコアのところをやらないといけないかな，という話ですね。グループダイナミックスというのは今言ったように，ミクロとマクロのダイナミックスというのが大事なんですよ，という話なんですけど，グループダイナミックス学会を僕が辞めた理由は，グループダイナミックス学会がダイナミックスを研究しているようには到底思えなかったからです。だってもう，「力学」と言っておきながら，その「力」だとか「測定」とかの定義も何もあったもんじゃなかったので。僕はそれでもグループダイナミックスをやろうと思っていたんですけど，一緒にやってくれる人があまりいなかったし，社会心理学とどこが違うんだろう，というのが明確にわからなかったので辞めた，というところですね。

ところで，マクロレベルの社会心理学は現象を扱うことが多いですね，という話をしましたけれども，社会心理学の研究が社会的な刺激に対する個人の反応，というようなところに落ち込んでたら，それだけの話なの？というのも気になります。オールポートは多分，もう少し広いことを考えてたんじゃないかと思います。そんなことも含めて，もう少し十分社会科学的なものとしての，社会心理学が捉えるべき何か，みたいなのをしっかりとこの授業を通じて書いていければなと思っております。

みんな疲れてきた？大丈夫？ちょっと休憩入れたりします？それとももう1回目からこんなに喋ると思ってなかったとかですか？何かあったらSli.doになにか書いてね。
僕は関西の出身の，大阪で25年生きていた人間ですから，授業をやるときに言うんですけど，皆さんがシーンと黙っておられるとですね，スベってると思うんです。スベっていると思うとどうなるかというと，もっと喋ってしまうんですよ。だから，早く黙れと思っているのであれば，何か言っていただければ少なくとも一人相撲みたいになっていないと思うのでいけるんですけど。なんかだんだんずっとスベり続けている気がしてきた。あと，ちょっとだけ社会心理学がどういう研究をしてきたかということと，この授業で何をしていくかということだけ，ちょっとだけ話すので，もう少しだけ我慢してください。まあまあ，喋らないといけないことにメモが入ってきたので，それぐらいは言っておこうと思います。

\section{社会心理学の三大主義}

僕もそんなこんなでずっと社会心理学って何学なんだろうと思いながらやってきておりますから，これがきっと社会心理学のコアなんだろうと言われる，社会心理学の三大主義，三つの基本的な考え方，ということを説明しておきます。一応ね，これでも社会心理学の授業をするときのテンプレみたいなのがあって，それについて一通りは説明できるんです。それを書いておきますね(板書)。

社会心理学っていうのはこの三つが入っている学問ですよ，って言われたりします。
一つが論理実証主義，二つ目は方法論的個人主義，三つ目は状況主義，というやつです。

論理実証主義っていうのは，もう今時皆さん聞かないんですか？ そんなことない？ 僕が大学院生になったときは，とある社会学の先生から，僕は論理実証主義者だけど，君はどうかね？ みたいなことを聞かれて，変なことを聞く人だなと思ったんですけど，論理実証主義は，ロジックと証拠で持って話を進めていきましょう，という，いわゆるサイエンスの枠組みのところですよね。心理学が人文科学だとか，哲学みたいなところから独立してサイエンスになろうとしたというときに，当時の哲学的なウィーン学団とかが言っていた，論理実証主義という哲学的な考え方をベースに，科学というのはできるんだ，というふうにしていたわけですね。

論理実証主義は，ウィトゲンシュタインの論理哲学論の話をもとに，「こうやってサイエンスをしたらいいんだ」と元気づけられた人たちが支えている考え方なんですが，当のウィトゲンシュタインは，後期になった時に，後期ウィトゲンシュタインになった時に，「言葉は世界をそんなに全部反映させてないよ」と言って撤回しちゃってます。その中で，前期ウィトゲンシュタインの発想に基づいて作られているサイエンスを今でもやりましょう，というのは，ちょっと時代遅れ，かなり時代遅れかな，というふうにも思ったりするんですが，一応，哲学的なところからは足を洗ってデータを取って，仮説を立てて検証して進みましょう，というサイエンスのスタイルを取っている，というのがこれです。

二つ目が方法論的個人主義というものです。これは，社会学と社会心理学はここでもって分かれる，と言ってもいいかもしれません。方法論的個人主義というのは，何かと言いますと，データを取る単位が個人だよってことです。

社会学，なかでも計量社会学とかは確実ですが，データを集めますよね。だけど，社会学で有名なデュルケームの「自殺論」\parencite{SPjisaturon}が代表的な例かもしれませんが，社会学なんかは割と集計されたデータを使います。統計庁が出すような，例えば日本の20代の未婚率がどれぐらい，とか，集計されたデータみたいなのを元に議論するという側面があるわけです。が，社会心理学はデータを個人レベル，一人一人から取ろうとします。なぜかはわかりませんけど，一人一人から取ったデータを扱おうとします。一応，この辺が心理学なのかな，と個人的には思ったりします。

三つ目が，状況主義というやつです。これは社会心理学が他の心理学と違う側面，というところかなと思います。他の心理学が，例えば人格心理学，性格心理学，臨床心理学みたいなのが，個人の人間としての個体の性質を研究対象にしているのに対し，社会心理学は「状況が整えば，それに応じて個人の行動は変わりますよ」，あるいは「個人は状況の方にこそ大きく影響されていますよ」というような考え方を持っている。なので，その人の思想，心情，心持ちがどうであろうと，状況で大分コントロールできるんじゃないですか，というふうな，個人以外の方に注目しましょう，というのが社会心理学の主なものの見方ではあります。

なので，こういった観点を持っていれば，社会心理学っぽいというふうに言われるわけですね。方法論としては，さっきも言いました調査法，実験法，観察法というのを使います。特に，調査や実験のデータの取り扱いに関しては，個人に注目するわけではなくて，全体像みたいなのを扱うことが多いので，より統計的な手法が強く扱われますし，個人の実験デザインを検証するというよりは，複雑に絡み合った多変量を一気に解析しましょうというようなことをやるようになっている，というのが社会心理学の特徴ではないかと思います。

皆さんに色々チェックして挙げてもらった論文も，興味があったら読んでいただければと思うんですが，一応データを取って考えるだろうし，個人レベルからそのデータを取っているだろうし，個人の属性というよりは，状況，環境，他の変数による説明，機械的な説明を重視しているのではないかと思います。

さて。でも，です。ここで言っているのは方法論の話であって，方法論的に通過するところがあるかといって，それを学問の定義にできるのか，というのはやっぱり違うような気がしませんか。この方法を使っていれば，この考え方であれば，社会心理学です，というのは，学問全体を通じての総論というか，ゴールになる部分ではないです。

例えば，生物学はと言われたら，この生物，生命とは何かみたいな全体を通じたテーマがあるはずです。経済学とかというのは，経済という側面で世界，社会を切り取って，経済というのはこういう仕組みになっている。これをマクロなレベルか，ミクロなレベルで考えましょうという感じ。まず，学問領域全体としての統一的な問題意識があって，総論とか原論があって，それに向けて各論があるはずです。

なので，社会心理学みたいに，それはそのゴールは分からないけれど，こういったことをやっていて，これに人がたくさん関わっていて，いろんな分類の仕方もありますけど，何か，複数人集まったら社会だし，集まらなくても社会と言っているわけですから，一人以上の人間が関わるのは全部社会心理学とかと言ってしまうと，社会心理学の境界が消えてしまうわけです。僕はそこを懸念しているし，それを踏まえた上で，我々が考えるべき社会心理学とは何か，という話を，この授業で追っかけていきたいと思っているわけです。そのためのヒントとして，どういった考え方があるかというと，一つは社会心理学の歴史だと思っています。これは多分テキストにも指定したと思うんですが，この授業は前半は，この「アナトミア社会心理学」というのをベースに話をしていこうと思っています\parencite{anatomia}。

これは2002年に出ていました。21世紀になって出た本ですけど，要は社会心理学が50年くらい経った時に，振り返って社会心理学とはいかがなる学問であったかという，この授業のメインテーマを20年くらい前にゴリゴリと書いた本です。この本全体が言わんとしていることから，時代は20年くらい進んでいますから，その辺のお話も含めてお伝えできればな，と思っています。次回は，この辺を中心に，社会心理学における黎明期から始めて，今はどういう社会心理学を考えて実践するようになってきているか，という話をしていこうと思っています。一つのアプローチの仕方として，社会心理学の歴史的なものを，この授業では扱う，ということです。

二つ目は，社会心理学だけではなくて，学問の背後にある思想的，科学観をテーマにしたお話を次にしていきたいと思っています。哲学と言うと，そのような気もしますが，そうではなくて，ここでお話したいのは「システム論」というものです。

システム論って皆さん聞いたことがあるかどうか分かりませんが，英語でSystemというのは，日本語で言うとこの「系」の一文字で書かれるものです。20年ぐらい前には，私が大学院生になった時には，複雑系という言葉がものすごく流行りまして，「これからは複雑系科学だ」と言われたんですけれども，最近はもうそんなことが言われなくなってきた。あるいは当然になってきたかもしれませんが，その「系」の話をしたいと思います。このシステムに関する学問というのは，実は近代科学ができてきた時に，生命科学分野におけるvon Bertalanffyという人の「一般生命システム論」というやつから話が始まって，化学の世界で「自己組織化系」という話に引き，その次に「自己組織化系」，「オートポイエーシス」という話に進んでいます。

もし皆さんが，「今コスギは一体何を言ってるんだ」と思ったのであれば，成功です。何かというと，皆さんが聞いたことのないような話をしたいので。システムの考え方の根拠みたいな話をしたいと思います。このシステム論の話は，今言った「オートポイエーシス」の次に「複雑系」，そして複雑系の次に「ネットワーク科学」というものに合流します。そういった考え方の流れを追いかけていきたいと思います。このシステム論というのは，僕はこれからの，あるいはこれまでもそうだったのかもしれませんが，社会心理学の中心的な話なんではないかと考えています。システムというのは，始まったらちゃんと説明しますが，要は全体のことを表しています。いろんな要素が複合的に組み合わさって出来上がった全体のことをシステムというわけですが，個人と社会というのは，まさにシステムのはずなわけですから，これをどういうふうなものとして捉えるのか，ということについてのお話をやっていきたいなと思います。

こういうのもね，ちょっとふわふわしていて，個々の研究に直接関わる話ではないのですが，ものの見方に直結する話なので，その辺のお話ができればと思っています。

で，この授業の最後，後半の方には，社会心理学の研究で行われている方法論的な特徴のところを切り出して，「今やっている方法論で良いのでしょうか？」という話をしたいと思います。この辺は，私が今の大学で統計屋さんという側面を持っているからですけれども。そもそも僕は社会心理学と統計の両方に，最初から二足のわらじを履いていたわけです。その二足のわらじを履いていた人間からいくと，社会心理学で行われている研究のアプローチだとか，測定していることというのは何なのか，というところに疑問がありますので，その辺についての統計的なお話を，パート3という感じで，最後の2，3回でお話ができればと思っています。

そんな感じで，この授業を進めていきたいと思います。適当に黒板に字しか書かないような感じで，オールドファッショナルなスタイルでやりたいんですけれども，皆さんにとっては聞いているだけだと，しんどい話かもしれません。でもこんなふうにして，皆さんと一緒に「社会心理学って何なのか？」というようなことを，個々の研究にこだわったり掘り下げたりするというよりは，メタ的な観点でお話しできればなと思っています。

だらだらと一方的に喋り続けて申し訳ありませんでした。Sli.doを使って，皆さんにコメントをいただければと思っているんですが，授業が終わったら，一応僕の話を聞いて，「こんなふうに考えましたよ」とか，「これはよかった，ここは悪かった」みたいな話があれば。できればフィードバックが欲しいので，LMSの方にアンケートも出しておきましたから，よかったら答えてください。

何もなかったら，僕は今日は100分くらいスベり続けたんだと反省して，来週も頑張ってスベり続けたいと思います。

それでは，初回から長々と失礼いたしました。今日の授業はここまでにしておきます。
% 引用文献
\printbibliography[title=引用文献]
\end{document}
